\documentclass[12pt]{article}
\usepackage[utf8]{inputenc}
\usepackage[T1]{fontenc}
\usepackage{hyperref}
\title{Nyankpa Language Resources for K-12 Curriculum Development}
\author{Compiled by Claude Sonnet 5.5 (claude-sonnet-5-5) for Andr\'as Kornai}
\date{}
\begin{document}
\maketitle

\section{Introduction}
Nyankpa (also Yeskwa, Nyenkpa, Yasgua) is a Plateau (Niger-Congo) language of central Nigeria. It is spoken mainly in Jema'a LGA of Kaduna State and Karu LGA of Nassarawa State. Search results give about 70,000 speakers (2012). Four varieties are named: Ejung, Ejira, Engbe and Ogbom. Wikipedia names Tattara as the prestige dialect and the basis of the standard form. Hausa is the wider language of communication.

\section{Primary Educational Resources}
Joshua Project reports that a Nyenkpa Language Project began in 2009 and published literacy primers. The search results did not show the primers themselves, so no URL for them can be given here.

\section{Linguistic and Reference Materials}
\begin{itemize}
\item Wikipedia, Nyankpa language: \url{https://en.wikipedia.org/wiki/Nyankpa_language}
\item Blench, 2008, listed in Glottolog: \url{https://glottolog.org/resource/reference/id/25430}
\item Webonary page for Nyankpa, probably a dictionary site (I did not check its content): \url{https://www.webonary.work/nyankpa/?p=9444}
\end{itemize}

\section{Religious and Community Resources}
\begin{itemize}
\item Scripture Earth index for Nyankpa: \url{https://www.scriptureearth.org/00i-Scripture_Index.php?iso=yes}
\item Joshua Project people group record: \url{https://joshuaproject.net/people_groups/16029}
\end{itemize}

\section{Digital and Audio Resources}
Global Recordings Network has audio and slideshow materials: \url{https://globalrecordings.net/en/language/yes}

\section{Curriculum Development Context}
Some literacy and scripture materials exist. The page contents were not verified beyond the search summaries. Whether Nyankpa is taught in schools was not found.
\end{document}